
\chapter{The faces of data}
\label{sect:faces}
The main purpose of this chapter is to raise awareness of the different ways
data are used, and how and why data should be structured for a specific purpose.
The reader is also acquainted with a few common terms related to data
management.

%%%%%%%%%%%%%%%%%%%%%%%%%%%%%%%%%%%%%%%%%%%%%%%%%%%%%%%%%%%%%%%%%%%%%%%%%%%%%%%
\section{Introduction}
\label{sect:faces:introduction}
Data are created and used for many different purposes, and it should therefore
not be surprising that the same data occurs in different guises that are
geared for a specific type of use.  One of the most common mistakes made in
data analyses or data management is to take data that is structured for one
purpose, say presentation, and try and use it in that form for another purpose like
statistical analyses. The purpose of this chapter is to give an overview
of the main types of uses of data, the way data is typically structured for
each purpose, and to compare their relative advantages and disadvantages.




The main uses of data distinguished here are storage, transport, analyses, and
presentation. Each use comes with its own conventions of how to
structure and manipulate data that stem from the different demands that the
different uses imply. For example, for reliable storage, retrieval and updating
of information it is beneficial to store each piece of information exactly
once, while it may be beneficial for interactive data analyses to have multiple
copies of the same data available.  Similarly, when transporting data from one
computer system to another, it is important that data can be `packed' and
`unpacked' quickly with a variety of tools. When presenting data, it is
important that the relevant information can be grasped quickly by users.


Working with data is more comfortable when data is structured especially for a
specific purpose. As a trade-off the other purposes suffer: data formatted for
communicating between web services is hardly fit for human interpretation. One
important aspect of data management is therefore the ability to recognize the
task at hand and transform data into a suitable format. 

As an example, consider a table of (fictional) CO$_2$ emission figures by fuel
source. Pay special attention to how the data are formatted.
\begin{equation}
\begin{tabular}{rl}
\multicolumn{2}{l}{CO$_2$ emission (kiloton per annum)}\\
\hline
\textbf{fuel}    & \textbf{emission}\\
\hline
\multicolumn{1}{l}{Petrol}      & 215 \\
-- of which bio                  & \phantom{0}75  \\
\multicolumn{1}{l}{Diesel}     & 456 \\
-- of which bio                  & \phantom{0}89  \\ 
\hline
\end{tabular}
\label{eq:present}
\end{equation}
This type of format is very common. One frequently encounters it for example in
official statistics publications or in financial reports of businesses.
Table~\eqref{eq:present} is optimized so people can quickly eyeball the
contribution of biofuels as a fraction of total emissions for each fuel type.
As a trade-off this table is not very suited for analysis. For example, to
derive the contribution of regular fuels, one needs to subtract the emission
value in the second line from the value in the first line, and do this again
for line three and four. Performing such calculations is much easier, and less
error-prone when the analyst does not have to worry about in what line which
data item is stored exactly.


This situation can be improved by realizing that the first column actually
hides a two-level classification: fuels are split into Petrol and Diesel
while each of those is split into the bio and regular type. Knowing this, the
same information can be represented as follows.
%
\begin{equation}
\begin{tabular}{rrl}
\hline
\textbf{fuel}      & \textbf{type}  & \textbf{emission}\\
\hline
Petrol & regular  & 140 \\
Petrol & bio      & \phantom{0}75\\
Diesel& regular  & 367 \\
Diesel& bio      & \phantom{0}89\\
\hline
\end{tabular}
\label{eq:analyse}
\end{equation}
%
The totals represented in the first table are no longer explicitly represented in
Table~\eqref{eq:analyse}. This is because the totals in Table~\eqref{eq:present}
present the result of an analyses while Table~\eqref{eq:analyse} represents the
basis of an analyses, that is: the raw underlying data. The benefit of the
second representation over the first is that it allows for more general analyses
like computing the total contributions of regular and biofuels.

A good way to think of the difference is that Table \eqref{eq:present} is a
presentation of results, while Table \eqref{eq:analyse} represents a data set
for analyses.  The advantage of the second for over the first is that it allows
general analyses, like computing the total emission contribution of regular
fuels.  Thus, although both tables communicate the same information they differ
in area of application. Unfortunately, the presentation form is very often used
in analyses, especially in spreadsheet software.


Next, have a look at the following format, where the data has been transformed
to a machine-readable string of characters.
%
\begin{equation}
\begin{minipage}{0.9\linewidth}
\begin{Verbatim}[breaklines=true]
[{"fuel":"petrol","type":"regular","emission":140}, {"fuel":"petrol","type":"bio","emission":75}, {"fuel":"diesel","type":"regular","emission":367}, {"fuel":"diesel","type":"bio","emission":89}]
\end{Verbatim}
\end{minipage}%
\label{eq:machine}
\end{equation}
%
Here the little $\hookrightarrow$ indicates that a line break was inserted only
to fit the code on the page. The actual string has no spaces and no line
endings or newlines.  The format presented here is called javascript object
notation (JSON).  The main purpose of this format is to provide a general way
to represent all sorts of data in a standardized way so it can be sent from one
piece of software to another. Now there are many of such formats, but this
particular one is very commonly used for communication between web services and
clients. Clearly, this format is not meant for presentation to people or for
statistical analyses. One advantage of JSON is that it is very generic: data
need not fit in a table to be stored in a JSON string, and in principle any
kind of data can be packed as JSON. Another advantage is that all wide-spread
programming languages have libraries that allow programmers to read, write, and
manipulate JSON data. It is thus a good format to interchange data between
different systems.


Finally, consider the same information but now organized as it would be, 
when stored in a database.
%
\begin{equation}
\begin{tabular}{ll}
\multicolumn{2}{l}{Fuel}\\
\hline
\textbf{id} & \textbf{name}\\
\hline
11 & petrol  \\
12 & diesel \\
\hline
\end{tabular}\textrm{, }\:
%
\begin{tabular}{ll}
\multicolumn{2}{l}{Type}\\
\hline
\textbf{id} & \textbf{name}\\
\hline
1 & regular \\
2 & bio \\
\hline
\end{tabular}\textrm{, }\:
%
\begin{tabular}{llll}
\multicolumn{3}{l}{Emission}\\
\hline
\textbf{id} & \textbf{fuel} & \textbf{type} & \textbf{amount} \\
\hline
120 & 11 &1 &140\\
121 & 11 &2 &\phantom{0}75\\
123 & 12 &1 &367\\
124 & 12 &2 &\phantom{0}89\\
\hline
\end{tabular}
\label{eq:database}
\end{equation}
Now each classification and each concept is given it's own table. There is a
`Fuel' table where each fuel is given a meaningless identifying number, and a
name. The same happened for the subclass `Type'. The table on emission sources
now has a unique id for each emission, and the fuel and type are identified by
the id's in their respective table. The reason data is stored in such a format
in database systems is that they are optimized for consistency. For example, if
one needs to replace the fuel name `petrol' with `gasoline', then the formats
presented in \eqref{eq:present}, \eqref{eq:analyse}, and \eqref{eq:machine}
require multiple changes in multiple places, which is both inefficient and
error-prone.  The database format ensures that in case of an update, only a
single item needs to be touched and everything else is updated immediately.
This benefit of easy and consistent updates comes with a cost: the data in the
Emission table is not understandable unless one knows the other tables as well.
We shall go into more detail when we discuss relational data in
Chapter~\ref{sect:relational}.

The above examples give some intuition to why it is important to structure data
in a way that is suited for purpose. Below the main types of uses, specific
demands that come with it and typical tools are summarized.

\paragraph{Storage.} When data is formatted for storage, this is typically done
with the four basic operations in mind that are often referred as the CRUD
operations.
\begin{terminology}{CRUD}
Creating, Retrieving, Updating, and Deleting data.
\end{terminology}
To facilitate this, data storage systems must have a way of uniquely
identifying each data item entrusted to them, regardless for example of the
order in which the data is stored.  The database example \eqref{eq:database}
satisfies extra demands such as a form of consistency, but there are other
storage solutions that do not satisfy this consistency but that do support CRUD
operations. Software for data storage is often characterised as either
relational or non-relational. The table set in \eqref{eq:database} is an
example of how data is structured in a relational database. Examples of
non-relational databases include key-value stores, document databases or graph
(network) databases. 

\paragraph{Transport.} When transporting data from one machine to another, it
is important that every value can be packed, send, unpacked and recognized
correctly. Over time many standards have been developed. Important formats
include the JSON format that was mentioned earlier and XML (eXtensible Markup
Language). Both formats are supported by many programming languages and they
can therefore be used by developers to exchange information between systems. 
The formats still leave a lot of choice for how to store data precisely. For 
example, in \eqref{eq:machine}, the order of variables in each record is
fuel--type--emission. Changing this order would still yield a valid JSON
format. For this reason, developers usually fix the way a JSON or XML document
must look like in something called a schema.
%
\begin{terminology}{Schema}
A schema describes the structure of a data format for a
specific type of data, within a chosen standard (e.g. XML, JSON).
\end{terminology}
%
Both \hlink{https://json-schema.org/}{JSON} and
\hlink{https://www.w3schools.com/xml/schema\_intro.asp}{XML} have their own
ways of denoting a schema. Some schemas are fairly simple, such as the JSON
schema defined by the European Statistical System for exchanging information
about
\hlink{https://ec.europa.eu/eurostat/cros/content/validation-report-structure\_en}{data
validation results}. Other schemas are large and elaborate, such as the
\hlink{https://sdmx.org}{SDMX standard}, for standardized exchange of
statistical information between institutions. 


\paragraph{Analysis.} To analyze data interactively, it is important that it is
formatted so that typical operations such as filtering, plotting, and modeling
are made as easy as possible. The most important demands are that a data set
represents a single type of entity and that each data item represents the
properties of one unique entity. Furthermore, there should be no duplicate
entries and all properties should be stored in the correct type (integer,
numeric, nominal, ordinal, text).  
%
\begin{terminology}{Entity type, population, entity, population unit, attribute, variable}
In computer science and specifically in the area of databases, the term \textit{entity
type} refers to a collection of real-world objects or events such as persons, clicks
on a web page, buildings in some city, and so on. In statistics this concept is usually
referred to as a \textit{population}, or target population.

Similarly, in computer science an \textit{entity} is an instance of an entity
type, while in statistics one often refers to a population element as a
\textit{population unit}.

All elements of a population (entities of an entity type) share a set of
properties.  Computer scientists often refer to a property as
\textit{attributes} while statisticians call them \textit{variables} or
\textit{stochastic variables} when probabilistic properties of the data are
important.
\end{terminology}
%
In many, but not all cases (See Section~\ref{sect:faces:types}) this means that
data is represented in a tabular format such as \eqref{eq:analyse} where each
row represents one entity, each column represents one attribute and  each
entity appears only once. A good rule of thumb to decide whether a tabular data
set is fit for analyses is to answer the question: are summary statistics of
the attribute columns meaningful? Compare for instance
Tables~\eqref{eq:present} and~\eqref{eq:analyse}: is it meaningful to compute
the sum of the numeric column? 

\paragraph{Presentation.} The main purpose of data in a presentable format is
to convey a message as clearly and objectively as possible. This means that
presentation is adapted to the expected audience and one may even choose
different presentations for different audiences. Typical tools for data
presentation include spreadsheet software, or business intelligence software
like Tableau. Almost all statistical software packages also offer some kind of
support for data presentation.


In practice, the lines between various uses of data are often blurred by
software tools. Spreadsheet software is optimized for presentation of data but
also offer limited analytical and data management features. Many commercial
databases offer some form of in-database analytical capabilities mixing storage
and analyses while users of spreadsheet software often confuse a spreadsheet
file with a data storage format. Although this blurring can seem convenient it
is almost always better to separate concerns and choose a specialized tool and
data format for the task at hand. The main reason for this is that
reproducibility, maintainability and understandability suffer when a tool is
used for purposes for which it is not designed.  Notorious examples include
complex interlinked sets of spreadsheets where it is very difficult to track
how the results are derived. The lack of systematic testability as well as
their opaqueness renders spreadsheets all but useless for reliable and
reproducible research. They do shine when it comes to presentation, with
features like conditional formatting or interactive pivoting. A second example
of such misuse of a tool is when programming a data analyses in Structured Query
Language (SQL). SQL is optimized for CRUD operations on data bases, but
expressing analytical operations such as computin summary statistics or
model parameters cumbersome when expressed in pure SQL. This then tends to lead
to code that is hard to read and difficult to maintain. So as a general rule,
it is a good idea to choose the data structure and associated tools that are
designed for a specific type of use. 





%%%%%%%%%%%%%%%%%%%%%%%%%%%%%%%%%%%%%%%%%%%%%%%%%%%%%%%%%%%%%%%%%%%%%%%%%%%%%%%
\section{Data types and data manipulation}
\label{sect:faces:types}
We have seen that data can be formatted in different ways for different
purposes.  The structure we use ---a database schema, a JSON string, in principle
says nothing about the \emph{intrinsic} structure of data, which is the
discussion of the following paragraphs.

\paragraph{Relational data.} 
The running example of the previous section is a data set where there is little
intrinsic relation between the different data points. The only intrinsic
structure is the fact that each entity is endowed with a number of attributes.
The way this data is represented has no real meaning: the order of columns and
rows in Table~\ref{eq:analyse} is of no importance. As long as we know which
values belong to which attribute of which entity, the content of the data stays
the same. Any analyses performed on this dataset will be insensitive to
permutations of rows or columns as long as we refer to the columns by their
name rather than position. Data satisfying this requirement is called
`relational data': each relation (row in a table) represents values of
attributes of a single population unit. Typical operations that one performs on
relational data include filtering or sorting rows, or adding or deleting
columns. These operations consume relational data and produce relational data.
The result of certain analyses can also be meaningfully stored in a relational
data set. For example grouped aggregation.

There are many cases where data points represented in a data set do possess
intrinsic interrelations, and this always requires specialized theoretical and
computational tools to analyze. Let us look at some important practical
examples.

\paragraph{Time series.} A time series is a sequence of values (or
value-tuples) that correspond to a sequence of equidistant times or periods.
Examples include monthly GDP estimates, daily closing levels of the stock
market or quarterly estimates of immigration. Many outputs in official
statistics are as time series,  and much effort is spent on constructing them.
Important operations performed on economic time series include inflation
correction, taking account of calendar effects, decomposition into trend and
seasonal effects, decomposition into long-term trend and business cycle
effects, increasing or decreasing the frequency by disaggregation or
aggregation, and more. All these operations have in common that they they both
accept and produce time series. They make explicit use of the fact that data is
organized as a time series and are therefore not invariant under permutation of
the elements of a time series.

\paragraph{Geospatial data.} This data type is used to visualise and analyse
phenomena where geographical location or geospatial patterns are important.
Some examples from official statistics where geospatial analyses plays a role
include land use, population density, biodiversity, and environmental
accounting. In geospatial data, each value is related to a point, line, or
region (a polygon) that is represented in a certain co-ordinate system.
Typical operations include co-ordinate transformations, geospatial smoothing,
and merging or splitting of regions. Again, permuting the values with respect
to their geospatial features (point, line, region) would make such operations
useless. Also observe that again, the `typical operations' mentioned here both
consume and produce geospatial data.


\paragraph{Network data.} A network consists of nodes, connected by links.
Nodes represent entities while links represent connections between them.  For
example a client--server relationship between companies, a family link between
persons, or a shared employer between persons.  At the time of writing,
official statisticians are just beginning to explore the uses of network data,
but it is likely that a more `systemic' approach to describing the economy and
society based on network science will play a role in the future of official
statistics. For networks, the structure \emph{is} the data. In network analyses
one is interested in summarizing or using network structure as a means to gain
insight in the phenomenon described by it. Typical operations on networks
include subnetwork selection or CRUD operations like adding or deleting nodes
or links. Again (and notice this is becoming a theme now), these operations
consume and produce networks.


\paragraph{Other.} There are many more data types that are useful but not
naturally represented in relational format. Examples include images, videos, or
event logs. All these data types come with their own `natural' transformations.

By now, attentive readers have probably detected a recurring theme across
these examples. Each data type comes with a set of `natural' operations that
leave the data type intact. This phenomenon is so important that mathematicians
and computer scientist have a name for it.
\begin{terminology}{Algebraic closure} Given
a data type $T$, and an operation that manipulates data of that type.  We say
that $T$ is \emph{algebraicly closed} under this operation when the result of
the operation is again of the same type $T$.
\end{terminology}

In practice, things are easier when you work in an algebraically closed system
of operations. For example. Given a relational data set, we can develop a
script where we first filter, then sort, and then merge the data set with
another data set. These operations are easy to perform in sequence because the
output of one operation can serve as input to the next. If we discover a bug
and find out we need to merge before filtering, the only thing we need to do is
swap the two operations. Indeed, working within an algebraically closed system
feels so natural that one can focus on \emph{what} of the data manipulation,
rather than at the \emph{how}.

There are other advantages that are facilitated by algebraic closure. For
example, most modern database systems are capable of re-ordering instructions
to increase speed of execution. For example it is faster to first filter
records and then sort them then the other way around (because there is less to
sort that way). Algebraic systems generally come with a set of re-ordering
rules that can be taken advantage of by software that executes it. Important
examples of algebraic systems are the various query languages that implement
CRUD operations for different data types. Examples include Structured Query
Language (SQL) for relational data, Cypher for network data and Spatial SQL for
geographic data.


The main takeaway of this section is that it is useful to distinguish between
\emph{data transformations} where data is manipulated but the type does not
change, and \emph{type transformations} where data manipulation yields data of
a new type. Third, we saw in \S\ref{sect:faces:introduction} there are format
transformations that change the technical layout of data. 

A typical step in a statistical production system starts with reading data from
a file or database. This is a format transformation, possibly followed by a
type transformation. For example when a JSON string is read into a relational
format. Next, depending on the task several type-conserving transformations
and/or type-transforming data manipulations are carried out after which the
data is again formatted for output.


\section*{Where are we?}
We have seen that data is a representation of information, and that data may be
formatted in several ways depending on its intended use. We have also seen that
data types can be distinguished by looking at their intrinsic structure.
Depending on intrinsic structure, a number of `natural' operations are
associated with data that leave its characteristic structure intact. It is
useful to distinguish between formatting, type-conserving, and type-changing
data transformations.



\newpage{}
\section{Exercises}
\label{sect:exfaces}
\begin{enumerate}
  \item Explain why Table~\eqref{eq:present} is unsuited for storage or transport of data.
  \item Acquaint yourself with storage solutions by reading through the Wikipedia
  entries for \hlink{https://en.wikipedia.org/wiki/Database}{Database} and
  \hlink{https://en.wikipedia.org/wiki/NoSQL}{NoSQL}.
  \item Find out what OLTP and OLAP means.
  \item Explain why or not the following data structures are suitable for analyses.
  \begin{enumerate}
      \item 
          \begin{tabular}{|r|lll|}
          \hline
                    & Alice & Bob & Carol \\
          \hline
          Shoe size & 38    & 43  & 41    \\
          Income    & 3300  & 2800& 4000  \\
          \hline
          \end{tabular}
      \item 
          \begin{tabular}{|r|llll|}
          \hline
            Income   &   0k -- 20k   &  20k -- 40k  &  40k -- 80k   & 80K+ \\
          \hline 
           Amsterdam & 20\%          &   40\%       &   35\%        & 5\%  \\
           Rotterdam & 30\%          &   30\%       &   38\%        & 2\%  \\
           Den Haag  & 25\%          &   35\%       &   30\%        &10\%  \\
          \hline
          \end{tabular}

      \item
          \begin{tabular}{|r|ll|}
          \hline
           Name & Age  &has job \\
          \hline
           Dave & 36   &No      \\
           Eve  & 5    &Yes     \\
          \hline
          \end{tabular}
      \item
\begin{tabular}{|r|ll|}
\hline
 Sector      & costs  &profit \\
\hline
 Retailers   & 50     &10     \\
 Wholesalers & 20     & 5     \\
 Total       & 70     &15     \\
\hline
\end{tabular}


  \end{enumerate}
\item Name three operations on image data that leave the data type intact.
\item Explain why each of the following statements is true or false.
\begin{enumerate}
  \item The natural numbers $\mathbb{N}=\{1,2,\ldots\}$ are closed under multiplication.
  \item The rational number s $\mathbb{Q}$ are closed under addition.
  \item The real numbers $\mathbb{R}$ are closed under division.
  \item We denote with $\Sigma^*$ the set of all finite strings that can be
    created with an alphabet $\Sigma$. For example if $\Sigma=\{a,b,\ldots,z\}$
    then $\Sigma^*$ is the set of all finite concatenations of letters from the
    Latin alphabet, plus the empty string (the string without characters).
    Statement: $\Sigma^*$ is closed under concatenation of strings (pasting one
    string behind the other).
 \item The function $\textrm{zero}$ returns $0$ for each value it is given. For
       Example $\textrm{zero}(\pi)=0$. Statement: the integers $\mathbb{Z}=\{\ldots, -2, -1, 0, 1, 2,\ldots\}$ are closed under $\textrm{zero}$.
 \item The natural numbers are closed under $\textrm{zero}$.
 \item (Bonus question) The function $a_n$ adds the number $1$, $n$ times to
    its input.  For example $a_3(2)=2+1+1+1=5$. Statement: the natural numbers are
    closed under $a_\infty$.
\end{enumerate}
\end{enumerate}

\hfill{}Check your answers on page~\pageref{sect:ansfaces}.

